\section{배경}\label{sec:background}

\subsection{전개형 게임과 기호}\label{sec:bg-notation}
전개형 게임은 $G=(\cN,\cH,\cZ,A,P,\sigma_c,u,\cI)$로 쓴다.
\begin{itemize}
  \item $\cN=\{1,2,c\}$는 플레이어 집합이고, $c$는 우연이다. $i$는 한 플레이어, $-i$는 그 상대와 우연을 뜻한다.
  \item $\cH$는 루트에서 시작하는 행동열(이력)의 유한 집합이다. $h\sqsubseteq h'$는 $h$가 $h'$의 앞부분이라는 뜻이다. $\cZ\subseteq\cH$는 종료 이력이다.
  \item $A(h)$는 비종료 이력 $h$에서 둘 수 있는 행동집합이고, $P(h)\in\cN$은 $h$에서 행동하는 플레이어다.
  \item $\sigma_c(h,\cdot)$는 우연 노드 $h$의 고정 분포이고, $u_i:\cZ\to\R$는 효용이다. 2인 제로섬이면 $u_1=-u_2$이고, $\Delta=\max_zu_1(z)-\min_zu_1(z)$를 효용 폭이라 한다.
  \item $\cI_i$는 플레이어 $i$의 정보집합들로, $P(h)=i$인 이력들의 분할이다. 같은 정보집합의 이력은 행동집합이 같다($A(I)$).
\end{itemize}

\paragraph{전략과 도달확률.}
행동 전략 $\sigma_i$는 각 $I\in\cI_i$에 분포 $\sigma_i(I)\in\Delta(A(I))$를 준다. $\sigma=(\sigma_1,\sigma_2)$를 프로필이라 한다.
\begin{itemize}
  \item $\pi^\sigma(h)$는 $\sigma$와 $\sigma_c$로 $h$에 이를 확률이다. 이것을 플레이어 $i$의 행동 확률의 곱 $\pi_i^\sigma(h)$와 나머지(상대와 우연)의 곱 $\pi_{-i}^\sigma(h)$로 나눈다. 그러면 $\pi^\sigma(h)=\pi_i^\sigma(h)\pi_{-i}^\sigma(h)$다.
  \item $h\sqsubseteq z$이면 $\pi^\sigma(h,z)$는 $h$에서 $z$까지 행동 확률의 곱이다.
  \item $\cZ_I=\{(h,z):h\in I,\ z\in\cZ,\ h\sqsubseteq z\}$로 둔다.
  \item 기대 효용은 $u_i(\sigma)=\sum_z\pi^\sigma(z)u_i(z)$다.
\end{itemize}

\paragraph{완전 기억.}
같은 정보집합 $I\in\cI_i$의 모든 이력에서 플레이어 $i$가 지나온 (정보집합, 행동)의 열이 같으면, 그 게임을 완전 기억 게임이라 한다.
이때 $\pi_i^\sigma(h)$는 $h\in I$에 대해 같으므로 $\pi_i^\sigma(I)$로 쓴다. 정보집합 $I$ 앞에 $i$가 둔 수의 개수를 $I$의 \emph{자기 깊이}라 한다.
완전 기억이면 행동 전략만으로 모든 혼합 전략을 표현할 수 있다~\citep{kuhn1953extensive}.
또 행동 전략 $\sigma_i$와 종료 이력마다의 자기 행동 확률 $x_i(z)=\pi_i^\sigma(z)$(실현 계획, sequence form~\citep{vonstengel1996efficient})가 일대일로 대응한다. $u_i$는 각 플레이어의 실현 계획에 대해 선형이다.

\paragraph{균형과 exploitability.}
$e_i(\sigma)=\max_{\sigma'_i}u_i(\sigma'_i,\sigma_{-i})-u_i(\sigma)$를 플레이어 $i$의 최선 응답 이득이라 한다.
$\NC(\sigma)=\sum_ie_i(\sigma)$가 0이면 $\sigma$는 내쉬 균형이다. $\NC(\sigma)\le\epsilon$이면 $\epsilon$-균형이라 한다.
이 논문은 실험에서 exploitability를 $\NC/2$로 보고한다.

\subsection{후회와 평균 전략}\label{sec:bg-regret}
반복 $t=1,\dots,T$의 프로필 $\sigma^t$에 대해, 플레이어 $i$의 전체 후회는 다음과 같다.
\[
R_i^T=\max_{\sigma'_i}\sum_{t=1}^T\Big(u_i(\sigma'_i,\sigma^t_{-i})-u_i(\sigma^t)\Big).
\]
평균 전략 $\bar\sigma^T$는 실현 계획의 평균 $\bar x_i=\frac1T\sum_tx_i^t$에 대응하는 행동 전략이다. 정보집합마다 쓰면 자기 도달확률로 가중한 평균이다.
\begin{equation}\label{eq:avg}
\bar\sigma_i^T(I,a)=\frac{\sum_{t=1}^T\pi_i^{\sigma^t}(I)\,\sigma^t(I,a)}{\sum_{t=1}^T\pi_i^{\sigma^t}(I)}.
\end{equation}
$\bar\sigma^T$는 반복 전략들을 균등하게 섞은 혼합 전략과 실현 동치다.

\begin{theorem}[평균 전략의 NashConv; folk theorem]\label{thm:folk}
2인 제로섬 게임에서 두 플레이어가 같은 프로필 열 $\sigma^1,\dots,\sigma^T$를 쓰면 $\NC(\bar\sigma^T)=(R_1^T+R_2^T)/T$다.
\end{theorem}
\begin{proof}
상대의 실현 계획에 대한 선형성으로 다음이 성립한다.
\[
u_i(\sigma'_i,\bar\sigma^T_{-i})=\frac1T\sum_tu_i(\sigma'_i,\sigma^t_{-i}),\qquad
\max_{\sigma'_i}u_i(\sigma'_i,\bar\sigma_{-i}^T)=\frac1T\Big(R_i^T+\sum_tu_i(\sigma^t)\Big).
\]
제로섬이므로 $\sum_iu_i(\bar\sigma^T)=0$이고 $\sum_iu_i(\sigma^t)=0$이다. 두 플레이어에 대해 더하면 된다.
\end{proof}

\subsection{CFR}\label{sec:bg-cfr}
\begin{definition}[반사실적 가치와 후회~\citep{zinkevich2007regret}]
$I\in\cI_i$의 반사실적 가치와 행동 $a$의 반사실적 가치를 다음과 같이 쓴다.
\[
v_i(\sigma,I)=\sum_{(h,z)\in\cZ_I}\pi^\sigma_{-i}(h)\,\pi^\sigma(h,z)\,u_i(z),\qquad
v_i(\sigma,I,a)=\sum_{(h,z)\in\cZ_I,\ ha\sqsubseteq z}\pi^\sigma_{-i}(h)\,\pi^\sigma(ha,z)\,u_i(z).
\]
정의상 $v_i(\sigma,I)=\sum_a\sigma(I,a)v_i(\sigma,I,a)$다.
프로필 $\sigma$에서의 순간 반사실적 후회는 $r^\sigma(I,a)=v_i(\sigma,I,a)-v_i(\sigma,I)$다.
반복 $t$의 순간 후회는 $r^t=r^{\sigma^t}$이고, 누적 후회는 $R^T(I,a)=\sum_{t\le T}r^t(I,a)$다.
\[
\psi(x)=\max\big(0,\max_ax_a\big)
\]
로 둔다. Zinkevich 외의 "직접 반사실적 후회의 양수 부분"은 $\psi(R^T(I))/T$다.
\end{definition}

CFR은 정보집합마다 후회 매칭으로 다음 전략을 정한다.
\begin{equation}\label{eq:rm}
\sigma^{T+1}(I,a)=\frac{R^T(I,a)^+}{\sum_bR^T(I,b)^+}\qquad(\text{분모가 0이면 균등 전략}).
\end{equation}
CFR의 정리는 세 부분으로 되어 있다.
\begin{enumerate}[label=(\roman*)]
  \item 전체 후회가 정보집합별 후회로 묶인다(정리~\ref{thm:cfr-decomp}).
  \item 정보집합마다 후회 매칭이 후회를 $O(\sqrt T)$로 묶는다(정리~\ref{thm:rm}).
  \item 이 둘과 정리~\ref{thm:folk}로 평균 전략이 균형에 수렴한다(따름정리~\ref{cor:cfr}).
\end{enumerate}
아래에서 (i)은 Zinkevich 외의 재귀적 증명 대신 성능 차이 항등식(보조정리~\ref{lem:perfdiff})으로 증명한다. 같은 항등식을 5절에서 다시 쓴다.

\begin{lemma}[성능 차이 항등식]\label{lem:perfdiff}
완전 기억 게임에서 임의의 프로필 $\sigma$와 플레이어 $i$의 임의의 순수 전략 $\beta$에 대해 다음이 성립한다.
\[
u_i(\beta,\sigma_{-i})-u_i(\sigma)=\sum_{I\in\cI_i}\pi_i^\beta(I)\,r^\sigma\big(I,\beta(I)\big).
\]
여기서 $\pi_i^\beta(I)\in\{0,1\}$은 $\beta$를 따를 때 $I$에 이르는지를 나타낸다.
\end{lemma}
\begin{proof}
$q(h,a)=\sum_{z\sqsupseteq ha}\pi^\sigma(ha,z)u_i(z)$로 둔다. $D$를 자기 깊이의 최댓값이라 하자.
$d=0,\dots,D+1$에 대해 $\sigma^{(d)}$를 "자기 깊이 $<d$인 정보집합에서는 $\beta$, 깊이 $\ge d$에서는 $\sigma_i$"로 둔다. 그러면 $\sigma^{(0)}=\sigma_i$, $\sigma^{(D+1)}=\beta$다.

$\sigma^{(d+1)}$과 $\sigma^{(d)}$는 깊이 $d$인 정보집합에서만 다르다.
\begin{itemize}
  \item 그런 $I$의 이력 $h$에는 두 전략 모두 확률 $\pi_i^\beta(I)\pi_{-i}^\sigma(h)$로 이른다. 그 앞의 자기 행동은 모두 $\beta$이기 때문이다.
  \item $h$ 뒤로는 $\sigma^{(d+1)}$이 $\beta(I)$를 둔 뒤 $\sigma$를 따르고, $\sigma^{(d)}$는 처음부터 $\sigma$를 따른다.
  \item 깊이 $d$의 정보집합을 지나지 않는 종료 이력의 확률은 둘이 같다.
\end{itemize}
따라서 다음이 성립한다.
\begin{align*}
u_i(\sigma^{(d+1)},\sigma_{-i})-u_i(\sigma^{(d)},\sigma_{-i})
&=\sum_{I:\,\text{깊이 }d}\pi_i^\beta(I)\sum_{h\in I}\pi^\sigma_{-i}(h)\Big(q(h,\beta(I))-\sum_b\sigma(I,b)q(h,b)\Big)\\
&=\sum_{I:\,\text{깊이 }d}\pi_i^\beta(I)\,r^\sigma(I,\beta(I)).
\end{align*}
$d=0,\dots,D$에 대해 더하면 왼쪽이 $u_i(\beta,\sigma_{-i})-u_i(\sigma)$로 소거된다.
\end{proof}

\begin{theorem}[후회 분해~\citep{zinkevich2007regret}]\label{thm:cfr-decomp}
완전 기억 게임에서 임의의 프로필 열에 대해 $R_i^T\le\sum_{I\in\cI_i}\psi\big(R^T(I)\big)$다.
\end{theorem}
\begin{proof}
$\sum_tu_i(\sigma'_i,\sigma^t_{-i})$는 $\sigma'_i$의 실현 계획에 대해 선형이다. 그래서 최댓값은 다면체의 꼭짓점, 즉 어떤 순수 전략 $\beta$에서 얻어진다.
보조정리~\ref{lem:perfdiff}를 $t$마다 쓰고 더하면 다음을 얻는다.
\[
\sum_t\big(u_i(\beta,\sigma^t_{-i})-u_i(\sigma^t)\big)=\sum_I\pi_i^\beta(I)R^T(I,\beta(I))\le\sum_I\psi(R^T(I)).
\]
마지막 부등식은 $\pi_i^\beta(I)\in\{0,1\}$과 $R^T(I,\beta(I))\le\psi(R^T(I))$에서 나온다.
\end{proof}

\begin{lemma}[순간 후회의 크기]\label{lem:rbound}
완전 기억 게임에서 $\sum_{h\in I}\pi^\sigma_{-i}(h)\le1$이다. 따라서 $\abs{r^\sigma(I,a)}\le\Delta_{u,i}=\max_zu_i(z)-\min_zu_i(z)$다.
\end{lemma}
\begin{proof}
$I$에 이르는 자기 행동열을 늘 두는 전략 $\beta$를 생각한다. $h\in I$에 대해 $\pi_i^\beta(h)=1$이고, 사건 "이력 $h$에 도달"들은 서로 배반이다.
그래서 $\sum_{h\in I}\pi_{-i}^\sigma(h)=\sum_{h\in I}\pi^{(\beta,\sigma_{-i})}(h)\le1$이다.
$q(h,a)-\sum_b\sigma(I,b)q(h,b)$는 효용 값들의 볼록결합의 차이이므로 그 절댓값은 $\Delta_{u,i}$ 이하다.
\end{proof}

\begin{theorem}[후회 매칭~\citep{hart2000simple,zinkevich2007regret}]\label{thm:rm}
식~\eqref{eq:rm}을 쓰면 모든 $I\in\cI_i$에서 $\psi(R^T(I))\le\Delta_{u,i}\sqrt{\abs{A(I)}\,T}$다.
\end{theorem}
\begin{proof}
$x^+=\max(x,0)$을 성분별로 쓴다.
\begin{itemize}
  \item 성분별로 $(R^{t-1}+r^t)^+\le\big((R^{t-1})^++r^t\big)^+$이므로 다음이 성립한다.
  \[
  \norm{(R^t)^+}_2^2\le\norm{(R^{t-1})^++r^t}_2^2=\norm{(R^{t-1})^+}_2^2+2\langle(R^{t-1})^+,r^t\rangle+\norm{r^t}_2^2.
  \]
  \item 가운데 항은 0이다. $(R^{t-1})^+$가 0이 아니면 $\sigma^t(I)\propto(R^{t-1})^+$이고, $\langle\sigma^t(I),r^t(I)\rangle=v_i(\sigma^t,I)-v_i(\sigma^t,I)=0$이기 때문이다.
  \item 귀납법과 보조정리~\ref{lem:rbound}로 $\norm{(R^T)^+}_2^2\le\sum_t\norm{r^t}_2^2\le\abs{A(I)}\,T\,\Delta_{u,i}^2$를 얻는다.
  \item $\psi(R^T)\le\norm{(R^T)^+}_\infty\le\norm{(R^T)^+}_2$다.
\end{itemize}
\end{proof}

\begin{corollary}[CFR의 수렴~\citep{zinkevich2007regret}]\label{cor:cfr}
$A_i=\max_{I\in\cI_i}\abs{A(I)}$로 두면 $R_i^T/T\le\Delta_{u,i}\abs{\cI_i}\sqrt{A_i}/\sqrt T$이다.
두 플레이어를 동시에 갱신하면 다음이 성립한다.
\[
\NC(\bar\sigma^T)\le\frac{\Delta_{u,1}\abs{\cI_1}\sqrt{A_1}+\Delta_{u,2}\abs{\cI_2}\sqrt{A_2}}{\sqrt T}.
\]
\end{corollary}

\paragraph{RM+와 CFR+.}
RM+는 $Q^t=(Q^{t-1}+r^t)^+$를 쌓고 $\sigma^{t+1}\propto Q^t$를 둔다~\citep{tammelin2014solving}. 다음 보조정리는 RM+의 후회가 어느 구간에서 재도 묶인다는 것이다. ASWISS가 분할 뒤에도 남은 bucket의 표를 초기화하지 않아도 되는 근거다.
\begin{lemma}[RM+의 구간 후회]\label{lem:rmplus}
$Q^0=0$이고 $\norm{r^u}_\infty\le\Delta$이면, 임의의 구간 $(s,t]$에 대해 성분별로 $\sum_{u\in(s,t]}r^u\le Q^t$이다. 또 $\norm{Q^t}_2^2\le\sum_{u\le t}\norm{r^u}_2^2$이다.
따라서 $\psi\big(\sum_{u\in(s,t]}r^u\big)\le\Delta\sqrt{\abs{A}\,t}$다.
\end{lemma}
\begin{proof}\mbox{}
\begin{itemize}
  \item $Q^{u}\ge Q^{u-1}+r^u$이므로 $\sum_{u\in(s,t]}r^u\le Q^t-Q^s\le Q^t$다.
  \item $\sigma^u\propto Q^{u-1}$이므로 $\langle Q^{u-1},r^u\rangle=0$이다. 그래서 $\norm{Q^u}_2^2\le\norm{Q^{u-1}+r^u}_2^2=\norm{Q^{u-1}}_2^2+\norm{r^u}_2^2$이다.
\end{itemize}
\end{proof}
CFR+는 RM+에 교대 갱신과 반복 번호에 비례하는 평균 가중치를 더한 것이다~\citep{tammelin2015solving}. 교대 갱신의 보장은 Burch 외~\citep{burch2019revisiting}가 정리했다.

\subsection{MCCFR}\label{sec:bg-mccfr}
MCCFR~\citep{lanctot2009monte}은 종료 이력의 덮개 $\cQ=\{Q_1,\dots,Q_r\}$($\bigcup_jQ_j=\cZ$)를 정하고, 매 반복 블록 하나 $Q_j$를 확률 $q_j>0$으로 뽑는다.
$q(z)=\sum_{j:\,z\in Q_j}q_j$는 $z$가 뽑힌 블록에 들어 있을 확률이다. 뽑힌 블록만으로 다음 표본 반사실적 가치를 계산한다.
\[
\tilde v_i(\sigma,I\mid j)=\sum_{(h,z)\in\cZ_I,\ z\in Q_j}\frac{1}{q(z)}\,\pi^\sigma_{-i}(h)\,\pi^\sigma(h,z)\,u_i(z).
\]
$\tilde v_i(\sigma,I,a\mid j)$도 $ha\sqsubseteq z$인 항만 남겨 같은 식으로 정의한다. 표본 후회는 $\tilde r(I,a)=\tilde v_i(\sigma,I,a\mid j)-\tilde v_i(\sigma,I\mid j)$이고, CFR과 같은 후회 매칭에 넣는다.

\begin{lemma}[불편성~\citep{lanctot2009monte}, Lemma 1]\label{lem:unbiased}
$\E_{j\sim q}\big[\tilde v_i(\sigma,I\mid j)\big]=v_i(\sigma,I)$이고 $\E_{j\sim q}\big[\tilde v_i(\sigma,I,a\mid j)\big]=v_i(\sigma,I,a)$다. 따라서 $\E[\tilde r(I,a)]=r^\sigma(I,a)$다.
\end{lemma}
\begin{proof}
합의 순서를 바꾸면 다음을 얻는다.
\begin{align*}
\E_j\big[\tilde v_i(\sigma,I\mid j)\big]&=\sum_jq_j\sum_{(h,z)\in\cZ_I,\ z\in Q_j}\frac{\pi^\sigma_{-i}(h)\pi^\sigma(h,z)u_i(z)}{q(z)}\\
&=\sum_{(h,z)\in\cZ_I}\frac{\pi^\sigma_{-i}(h)\pi^\sigma(h,z)u_i(z)}{q(z)}\sum_{j:\,z\in Q_j}q_j=v_i(\sigma,I).
\end{align*}
행동 $a$의 경우도 같다. 후회는 두 추정량의 차이이므로 기댓값의 선형성으로 불편하다.
\end{proof}

\paragraph{외부 샘플링.}
갱신할 플레이어(traverser) $i$를 정한다.
\begin{itemize}
  \item 우연 노드에서는 $\sigma_c$로, 상대 노드에서는 $\sigma_{-i}$로 행동을 하나씩 뽑는다. 상대와 우연의 순수 전략 $\tau$를 하나 뽑는 것과 같다.
  \item 자기 노드에서는 모든 행동을 펼친다.
  \item 블록은 $Q_\tau=\{z:\ z\text{의 상대·우연 행동이 }\tau\text{와 같다}\}$이고, $q(z)=\pi^\sigma_{-i}(z)$다.
\end{itemize}
$\pi^\sigma_{-i}(z)=\pi^\sigma_{-i}(h)\pi^\sigma_{-i}(h,z)$이고 $\pi^\sigma(h,z)=\pi_i^\sigma(h,z)\pi_{-i}^\sigma(h,z)$이므로 중요도 가중이 약분된다.
\[
\tilde v_i(\sigma,I\mid\tau)=\sum_{(h,z)\in\cZ_I\cap Q_\tau}\pi_i^\sigma(h,z)\,u_i(z).
\]
즉 뽑힌 행동을 따라가며 자기 행동만 자기 확률로 펼치면 된다.
완전 기억이면 같은 $I$의 두 이력은 처음 갈라지는 곳이 상대나 우연의 노드다. 그래서 한 순회에서 $I$의 이력은 많아야 하나만 $\tau$와 맞는다. 따라서 $\abs{\tilde r(I,a)}\le\Delta_{u,i}$다.

\begin{theorem}[외부 샘플링 MCCFR의 수렴 속도~\citep{farina2020stochastic}, 명제 2와 식 (6)]\label{thm:es}
두 플레이어가 외부 샘플링 추정량을 받는 CFR을 동시에 돌린다고 하자. $\tilde R_i^T$를 플레이어 $i$의 CFR이 추정된 손실에 대해 갖는 후회라 하자.
임의의 $p\in(0,1)$에 대해 확률 $1-2p$ 이상으로 다음이 성립한다.
\[
\NC(\bar\sigma^T)\le\frac{\tilde R_1^T+\tilde R_2^T}{T}+4\Delta\sqrt{\frac2T\log\frac1p}.
\]
외부 샘플링의 표본 손실은 상대와 우연을 $\tau$로 고정한 게임의 효용이다. 그래서 정리~\ref{thm:cfr-decomp}와 정리~\ref{thm:rm}이 그대로 적용되어 $\tilde R_i^T\le\Delta\abs{\cI_i}\sqrt{A_iT}$다. 따라서 확률 $1-2p$ 이상으로 다음이 성립한다.
\[
\NC(\bar\sigma^T)\le\frac{\Delta\big(\abs{\cI_1}\sqrt{A_1}+\abs{\cI_2}\sqrt{A_2}\big)+4\Delta\sqrt{2\log(1/p)}}{\sqrt T}.
\]
\end{theorem}
\begin{proof}[증명 개요]\mbox{}
\begin{itemize}
  \item 고정된 비교 전략 $\hat x$에 대해 $d_t=(\ell^t-\tilde\ell^t)^\top(x^t-\hat x)$는 마팅게일 차분이다($\ell^t$는 참 손실, $\tilde\ell^t$는 불편 추정). 크기는 $2\Delta$ 이하다.
  \item Azuma–Hoeffding 부등식으로 확률 $1-p$ 이상에서 $\sum_td_t\le2\Delta\sqrt{2T\log(1/p)}$다.
  \item 그래서 참 후회는 $\tilde R_i^T+2\Delta\sqrt{2T\log(1/p)}$ 이하다. 두 플레이어에 합집합 상계를 쓰고 정리~\ref{thm:folk}를 적용한다.
  \item 비교 전략 전체에 대한 처리 등 세부는 원문을 따른다.
\end{itemize}
\end{proof}

\begin{remark}
Lanctot 외~\citep{lanctot2009monte}의 원래 분석도 외부 샘플링 MCCFR의 평균 후회가 높은 확률로 $O(1/\sqrt T)$임을 보였다. 다만 실패 확률에 대한 의존이 $O(\sqrt{1/p})$였고, 정리~\ref{thm:es}는 이것을 $O(\sqrt{\log(1/p)})$로 줄인 것이다~\citep{farina2020stochastic}.
상수는 원문마다 표기가 다르므로 인용할 때 원문의 정의를 따른다.
\end{remark}

\subsection{정보 추상화와 불완전 기억}\label{sec:bg-abstraction}
정보 추상화 $\cB=(\cB_1,\cB_2)$는 각 플레이어의 정보집합 $\cI_i$를 bucket들로 나눈 분할이다.
\begin{itemize}
  \item 한 bucket $B$의 멤버들은 같은 플레이어의 것이고, 행동집합이 같다($A_B$). 행동의 대응은 미리 정한다.
  \item bucket 전략 $\sigma(B)$를 멤버에 그대로 쓰는 것을 lifting이라 하고 $\sigma(I)=\sigma(B(I))$로 쓴다. lifting된 전략은 원 게임 $G$의 행동 전략이므로 $G$에서 그대로 평가할 수 있다.
  \item bucket 위에서 CFR이나 MCCFR을 돌리면, bucket의 후회 갱신은 멤버들의 반사실적 후회의 합이다.
    \[
    r^\sigma(B,a)=\sum_{I\in B}r^\sigma(I,a).
    \]
\end{itemize}
서로 다른 자기 행동열을 거친 정보집합이 같은 bucket에 들면 그 추상화는 불완전 기억이다~\citep{waugh2009practical,lanctot2012no}.
그중 한 경로가 같은 bucket을 두 번 지나는 경우를 \emph{absent-mindedness}라 한다~\citep{piccione1997interpretation}.
\begin{itemize}
  \item 불완전 기억 추상 게임에서는 행동 전략이 모든 혼합 전략을 표현하지 못한다. 그래서 추상 게임에서 CFR의 보장이 일반적으로 성립하지 않는다~\citep{lanctot2012no}.
  \item 이 논문은 추상 게임의 균형 대신 원 게임 $G$의 NashConv를 직접 묶는다. 그래서 lifting된 전략이 $G$의 행동 전략이라는 사실만 쓴다(5절).
\end{itemize}

\subsection{CFR+IRA}\label{sec:bg-ira}
CFR+IRA~\citep{cermak2020automated}는 2인 제로섬 완전 기억 게임 $G$에서 거친 불완전 기억 추상화로 시작한다. 학습 중에 추상화를 세분하며 $G$의 균형으로 수렴한다.
원문 4.3절의 구성은 다음과 같다.

\paragraph{초기 추상화(원문 4.1.1절).}
플레이어 $i$의 정보집합 중 자기 행동열의 길이와 행동 수가 같은 것끼리 하나로 묶는다. 행동은 게임 기술의 순서대로 대응시킨다.

\paragraph{학습.}
추상화되지 않은 트리 전체를 매 반복 순회하며 CFR+를 돈다(교대 갱신, 정해진 지연 뒤부터 평균). 후회와 평균 전략은 추상 정보집합 단위로 저장한다.
멈춤 조건은 원 게임 최선 응답으로 잰 평균 전략의 exploitability가 $\epsilon$ 이하가 되는 것이다.

\paragraph{후회 상한 갱신(원문 4.3.1절, 수렴 보장용).}
\begin{itemize}
  \item 반복 열 $T_1<T_2<\cdots$를 $T_{j+1}-T_j=2^j$로 둔다.
  \item 각 $T_j$에서, 아직 추상 정보집합에 묶여 있는 원 정보집합 가운데 $k_b$개를 균등하게 뽑는다($\hat\cI_b$).
  \item 구간 $[T_j,T_{j+1})$ 동안 그 원 정보집합들의 후회를 따로 추적한다.
  \item 어떤 반복 $T'$에서 다음 조건이 성립하면, 그 $I$를 추상 정보집합에서 떼어 내고 후회와 평균을 새로 시작한다.
  \[
  R^{T'}_{T_j,\mathrm{imm}}(I)=\frac{\max_{a}\sum_{t=T_j}^{T'}\big(v_i(b^t_{I\to a},I)-v_i(b^t,I)\big)}{T'-T_j}>L_I^{T_j}(T'),\qquad L_I^{T_j}(t)\le\frac{\Delta_I\sqrt{\abs{A(I)}}}{\sqrt{t-T_j}}
  \]
\end{itemize}
즉 감시한 정보집합의 구간 평균 후회가, 그 정보집합이 따로 있었다면 후회 매칭이 보장했을 속도보다 느리게 줄면 떼어 낸다.

\paragraph{휴리스틱 갱신(원문 4.3.2절, 실전 속도용).}
\begin{itemize}
  \item 매 반복, 행동하는 플레이어의 정보집합 $k_h$개를 뽑아 \emph{그 반복 하나의} 후회 $r^h$를 기록한다.
  \item 각 정보집합에 $\gamma_I=\{a:\ r^h(a)\ge\max_{a'}r^h(a')-1/(5\sqrt t)\}$를 붙이고, $\gamma_I$가 같은 것끼리 나눈다.
  \item 가장 큰 묶음이 추적하지 않은 멤버를 데려간다. 새 묶음들의 후회는 0에서 시작한다.
\end{itemize}

\begin{theorem}[CFR+IRA의 수렴~\citep{cermak2020automated}, 정리 2]\label{thm:ira}
다음 식(원문 식 (26))의 $\alpha(\delta)$를 쓴다. \todo{원문 PDF에서 근호와 지수의 배치를 다시 확인할 것}
\[
\alpha(\delta)=\sqrt[3]{\,2^{\sqrt[\abs{\cI}]{\abs{\cI}!\,\log_{1-(k_b/\abs{\cI})^{\abs{\cI}}}(\delta)}\;+\;\abs{\cI}+3}}\;A_{\max}
\]
그러면 모든 $T\ge(\alpha(\delta)+1)^3$에서 확률 $1-\delta$로 다음이 성립한다.
\[
\frac{R_i^T}{T}\le\Delta\abs{\cI_i}\frac{\sqrt{A_{\max}}}{\sqrt T}+\Delta\abs{\cI_i}\frac{\alpha(\delta)}{\sqrt[3]T}+\Delta\abs{\cI_i}\frac1T\ \in\ O\!\left(\frac1{\sqrt[3]T}\right).
\]
\end{theorem}
\paragraph{증명의 구조(원문 보조정리 2–6).}
\begin{itemize}
  \item 처음 $T_{\lim}$ 반복 동안 추상화가 수렴을 막는다고 보고, 그 동안의 후회를 정보집합마다 최악 $\Delta T_{\lim}$로 센다.
  \item 그 뒤는 RM+의 $\Delta\sqrt{\abs{A(I)}T}$를 쓴다. 그래서 $R_i^T/T\le B_i(T)=\frac1T\big(\Delta\abs{\cI_i}\sqrt{A_{\max}}\sqrt T+\Delta\abs{\cI_i}T_{\lim}\big)$이다(원문 보조정리 3).
  \item 이 상한이 성립할 확률의 하한은 최악의 경우로 계산한다. 모든 원 정보집합을 정해진 순서로 하나씩 떼어 내야 하는 경우다.
    \[
    p_r(T,\bar R_i)=1-\Big(1-\big(k_b/\abs{\cI}\big)^{\abs{\cI}}\Big)^{\binom{\abs{\tau^T_{\bar R_i}}}{\abs{\cI}}}
    \]
    여기서 $\tau^T_{\bar R_i}$는 감시 구간 중 충분히 긴 것들의 열이다.
  \item $T_{\lim}=\lceil\alpha(\delta)T^{2/3}\rceil$로 두면 이 확률이 $1-\delta$ 이상이 된다. 그 결과가 $O(T^{-1/3})$다.
\end{itemize}

\paragraph{휴리스틱만으로는 수렴하지 않는다(원문의 반례).}
\begin{itemize}
  \item 한 추상 정보집합에 두 원 정보집합 $I',I''$가 묶여 있고, 행동 가치가 두 표 사이를 오간다고 하자(원문 표 1).
  \item 그러면 매 반복 두 정보집합의 최선 행동 집합 $\gamma$가 같아서 휴리스틱은 한 번도 나누지 않는다.
  \item 그러나 추상 정보집합은 균등 전략(기댓값 2.75)으로 수렴하는 반면, 따로 풀면 기댓값 5인 서로 다른 순수 전략이 균형이다.
  \item 한 반복의 순간 후회로 묶는 규칙은 진동하는 가치에서 이런 실패를 겪는다.
\end{itemize}
ASWISS의 라벨은 epoch 동안의 \emph{누적} 후회의 argmax다. 이 예에서 누적 후회는 두 정보집합에서 서로 다른 행동을 가리킨다(값의 크기가 1과 10으로 다르므로). 그래서 나뉜다.
일반적으로는 누적 라벨이 일치하면 숨은 후회가 0이다(보조정리~\ref{lem:label}).

\paragraph{한계.}
매 반복 추상화되지 않은 트리 전체를 순회한다. 정리~\ref{thm:ira}의 $\alpha(\delta)$는 정보집합 수에 지수적이라, 상한이 의미를 갖는 $T$가 현실적인 범위를 크게 넘는다.
원문 6.1절은 표본 기법과의 결합을 후속 과제로 남겼다.
